The incoming and the reflected wave form a standing wave; neighbouring minima (nodes) are $\lambda/2$ apart. Between the first and the last of \MiO\ minima lie $\MiO - 1 = 4$ gaps:
\begin{align}
 (n-1)\cdot\frac{\lambda}{2} = d \quad\Rightarrow\quad \lambda &= \lsF \\
   &= \frac{2\times\ds}{\MiO-1} \\
   &= \ls \approx \result{\lsP{-2}{2}}
\end{align}
The frequency:
\begin{align}
 f &= \fsF = \frac{\ncc}{\ls} = \fs \approx \result{\fsP{9}{2}}
\end{align}
So Heinrich Hertz measured the wavelength of his waves in 1888, and with the frequency of his circuit found their speed: that of light.
